\section*{Chapitre \ref{ch:qm:jump-processes} --- Processus à sauts}

\begin{solution}{exo:qm:jump-processes:1}
$e^{-5/12} = 65.9\%$~; par absence de mémoire, l'attente moyenne est de $1/5$ d'année, quoi qu'il se soit passé.
\end{solution}

\begin{solution}{exo:qm:jump-processes:2}
Moyenne $3 \times (-2\%) = -6\%$~; variance $3(0.02^2 + 0.04^2) = 0.006$, soit un écart-type de 7.75~\%.
\end{solution}

\begin{solution}{exo:qm:jump-processes:3}
Pour $s < t$, $N_t - N_s$ est indépendant de $\mathcal F_s$ de moyenne $\lambda(t - s)$, donc $\E_s[N_t - \lambda t] = N_s - \lambda s$~;
et $\E[(N_t - \lambda t)^2] = \Var(N_t) = \lambda t$.
\end{solution}

\begin{solution}{exo:qm:jump-processes:4}
$d(e^{X_t}) = e^{X_{t-}}(\sigma\,dW_t + \tfrac12\sigma^2dt) + e^{X_{t-}}(e^{\Delta X_t} - 1)$. Le terme de saut a pour compensateur
$e^{X_{t-}}\lambda(\E[e^Y] - 1)\,dt$, donc la dérive $b = -\tfrac12\sigma^2 - \lambda(\E[e^Y] - 1)$ ajoutée à $X$ fait de $e^X$ une
\omterm{def:qm:probability-at-speed:martingale}{martingale}.
\end{solution}

\begin{solution}{exo:qm:jump-processes:5}
$\psi(u) = 2(50/(50 - \iu u) - 1)$~; moyenne $\lambda\E[Y] = 2/50 = 0.04$, variance $\lambda\E[Y^2] = 2 \times 2/50^2 = 0.0016$.
\end{solution}

\begin{solution}{exo:qm:jump-processes:6}
$76.4/5 = 15.3$ à une semaine et $76.4/252 = 0.30$ à un an.
\end{solution}

\begin{solution}{exo:qm:jump-processes:7}
$(-0.1,\ 0.043,\ -0.00378,\ 0.001888)$ dans les deux cas~: les différences finies de la fonction
génératrice des \omterm{def:qm:jump-processes:cumulant}{cumulants} concordent avec les formules à environ $10^{-7}$.
\end{solution}

\begin{solution}{exo:qm:jump-processes:8}
Dans un modèle de Lévy, le kurtosis en excès décroît comme $1/t$~: 1.2 à 21 jours correspond à environ 25 à un jour.
Les rendements mensuels lissent précisément les sauts qui comptent pour le risque à un jour~; mesurer la queue à
l'horizon du risque.
\end{solution}

\begin{solution}{pb:qm:jump-processes:1}
\textbf{1.} $20.98 \times 0.98\% = 20.6\%$.
\textbf{2.} $\log_{10}\Phi(-20.98) = -97.3$.
\textbf{3.} $20\%/0.98\% = 20.4$ écarts-types~: $\log_{10}$ de la probabilité égal à $-92.2$.
\textbf{4.} $0.98\%\sqrt{252} = 15.6\%$.
\textbf{5.} Une volatilité gaussienne assez élevée pour rendre la journée plausible rendrait les journées ordinaires
plusieurs fois plus volatiles qu'elles ne le sont~; la queue et le corps ne peuvent pas tenir ensemble.
\textbf{6.} $\mu_J = -11.1\%$.
\textbf{7.} 13.0~\% par an (12.96~\%).
\textbf{8.} 31~\%~: $0.5/252 \times (0.111^2 + 0.05^2)$ sur $0.0098^2$.
\textbf{9.} $1 - e^{-0.5/252} = 0.20\%$.
\textbf{10.} Avec une probabilité de 0.116~\% par jour~: une fois tous les 3.4 ans.
\textbf{11.} Asymétrie $-4.64$, kurtosis en excès 76.4.
\textbf{12.} $76.4/21 = 3.64$.
\textbf{13.} Les \omterm{def:qm:jump-processes:cumulant}{cumulants} d'un \omterm{def:qm:jump-processes:levy}{processus de Lévy} sont linéaires en $t$, donc $\kappa_4/\kappa_2^2 \propto t/t^2$.
\textbf{14.} $-\lambda(\E[e^Y] - 1) = -0.5(e^{-0.111 + 0.00125} - 1) = +5.2\%$ par an, en plus de $-\tfrac12\sigma_c^2$.
\textbf{15.} 15.3 en théorie, 15.6 sur 400\,000 rendements simulés de cinq jours.
\textbf{16.} Un choix~: une telle journée dans l'historique ne permet pas d'estimer une fréquence~; le modèle encode un
jugement sur la fréquence à laquelle elle doit pouvoir se produire.
\textbf{17.} Regroupement de volatilité~: dans les données, le kurtosis décroît plus lentement que $1/t$ (chapitre~18).
\textbf{18.} Des options de vente à courte échéance, très hors de la monnaie, dont la valeur tient presque entièrement à la probabilité de saut (le
livre~5, chapitre~13).
\textbf{19.} Résultat à retenir~: \emph{vingt écarts types}~: la journée de 1987 fut une chute de 21 écarts types
avec une probabilité gaussienne de $10^{-97}$~; des sauts d'intensité 0.5 par an, de moyenne $-11.1\%$ et un écart-type
de 5~\%, portant 31~\% de la variance, font d'une chute de 20~\% un événement qu'on observe une fois tous les cinquante ans, au
prix d'un kurtosis en excès à un jour de 76 (3.6 à un mois).
\textbf{20.} Il modifie la forme de la loi à un jour, en plaçant de la masse dans la queue sans augmenter la
variance des journées ordinaires.
\end{solution}

\begin{solution}{iq:qm:jump-processes:1}
Que c'est le modèle, et non le marché, qui a tort~: sous celui-ci, le mouvement n'aurait jamais dû se produire. Vérifier les données
d'abord (une mauvaise cotation, une opération sur titres), puis les hypothèses du modèle (trajectoires continues, volatilité constante,
chocs gaussiens) et les mesures de risque qui en dépendent.
\iqlookfor{défaillance du modèle, pas la malchance.}
\end{solution}

\begin{solution}{iq:qm:jump-processes:2}
$1 - e^{-2}(1 + 2) = 59.4\%$~; une demi-seconde, quel que soit le temps écoulé depuis la dernière transaction.
\iqlookfor{probabilités de Poisson et absence de mémoire.}
\end{solution}

\begin{solution}{iq:qm:jump-processes:3}
Une couverture delta compense les petits mouvements au premier ordre~; un saut est un grand mouvement, et la valeur de l'option change
de toute la différence $V(S + \Delta S) - V(S)$, et non de $\Delta\cdot\Delta S$. Avec des tailles de sauts aléatoires, il existe une source de risque supplémentaire
que l'action ne permet de couvrir~: le marché est incomplet.
\iqlookfor{non-linéarité et incomplétude.}
\end{solution}

\begin{solution}{iq:qm:jump-processes:4}
$\psi(u) = \iu\gamma u - \tfrac12\sigma^2u^2 + \int(e^{\iu ux} - 1 - \iu ux\mathbf 1_{|x|<1})\,\nu(dx)$~: $\gamma$ une dérive, $\sigma^2$ la
variance brownienne, $\nu$ le nombre espéré de sauts par unité de temps selon leur taille.
\iqlookfor{les trois termes et le rôle de $\nu$.}
\end{solution}

\begin{solution}{iq:qm:jump-processes:5}
Les \omterm{def:qm:jump-processes:cumulant}{cumulants} varient proportionnellement à $t$, donc le kurtosis en excès décroît comme $1/t$ et un rendement mensuel est presque
gaussien. Dans les données, il décroît plus lentement, car la volatilité se regroupe et les accroissements ne sont pas
indépendants.
\iqlookfor{$\kappa_n(t) = t\kappa_n(1)$ et ce qui la met en défaut.}
\end{solution}

\begin{solution}{iq:qm:jump-processes:6}
Le prix actualisé doit être une \omterm{def:qm:probability-at-speed:martingale}{martingale}~: pour $S = S_0e^X$ avec $X = bt + \sigma W + J$, $b = r - \tfrac12\sigma^2 -
\lambda(\E[e^Y] - 1)$, le compensateur des sauts retranchant leur contribution moyenne. Sinon le modèle
admet un arbitrage contre le \omterm{def:qm:girsanov-and-changes-of-numeraire:numeraire}{compte monétaire}.
\iqlookfor{le compensateur dans la dérive.}
\end{solution}
